\chapter{Explicit Algorithmic State in \texttt{PlasticityRelation}}

\section{Why This Step Matters}

The previous chapter introduced \texttt{ConstitutiveIntegrator} as the
algorithmic injection point at the erased constitutive-handle boundary.  That
was necessary, but not sufficient.  As long as the nonlinear constitutive law
still owns its irreversible variables internally, the split remains only
partial.

This phase takes the first behavioural step beyond semantics:
\texttt{PlasticityRelation} can now be evaluated and committed against an
\emph{explicit external} algorithmic state.

\section{Old Situation}

Originally, \texttt{PlasticityRelation} stored a member
\(\alpha_n\) internally and all core operations were implicitly tied to that
member:

\begin{align}
\bsigma_{n+1} &= \texttt{compute\_response}(\bvarepsilon_{n+1}), \\
\mathbf{C}_{n+1}^{ep} &= \texttt{tangent}(\bvarepsilon_{n+1}), \\
\alpha_{n+1} &= \texttt{update}(\bvarepsilon_{n+1}).
\end{align}

That design is workable for a local prototype, but architecturally it creates
three problems:

\begin{itemize}
  \item the constitutive law is not clearly separated from the constitutive
        site;
  \item the local integration algorithm cannot be cleanly redirected to other
        state carriers;
  \item sharing or borrowing the law becomes semantically fragile because the
        relation still behaves as if it owned the site memory.
\end{itemize}

\section{New Explicit-State Contract}

\texttt{PlasticityRelation} now supports a second, more explicit interface:

\begin{lstlisting}[language=C++]
ConjugateT compute_response(const KinematicT& eps,
                            const InternalVariablesT& alpha) const;

TangentT tangent(const KinematicT& eps,
                 const InternalVariablesT& alpha) const;

void commit(InternalVariablesT& alpha,
            const KinematicT& eps) const;
\end{lstlisting}

Mathematically, the relation is now able to realize the return-mapping map

\begin{align}
(\bvarepsilon_{n+1}, \alpha_n)
\;\mapsto\;
(\bsigma_{n+1}, \mathbf{C}_{n+1}^{ep}, \alpha_{n+1})
\end{align}

without requiring \(\alpha_n\) to live inside the relation object itself.

\section{Compatibility Is Preserved}

The embedded-state API is still available:

\begin{itemize}
  \item \texttt{compute\_response(eps)}
  \item \texttt{tangent(eps)}
  \item \texttt{update(eps)}
  \item \texttt{internal\_state()}
\end{itemize}

Those methods now act as a compatibility layer around the explicit-state core.
Therefore:

\begin{itemize}
  \item existing tests and client code continue to work;
  \item the new architecture can progressively migrate site storage out of the
        relation;
  \item naked relations can still be used in the old stateful way when
        convenient.
\end{itemize}

\section{What Changed Internally}

The central change is not the public overloads by themselves, but the internal
factoring of the return-mapping kernel.  The predictor/corrector path now
receives \(\alpha_n\) explicitly:

\begin{itemize}
  \item elastic predictor uses the supplied plastic strain;
  \item effective trial stress uses the supplied hardening state;
  \item return-mapping builds \(\alpha_{n+1}\) from the supplied
        \(\alpha_n\), not from an implicitly owned member.
\end{itemize}

This makes the local nonlinear map much closer to its mathematical form.

\section{Interaction with \texttt{MaterialInstance}}

The typed constitutive site \texttt{MaterialInstance<Relation, ...>} now uses
the explicit-state path automatically when the relation satisfies
\texttt{ExternallyStateDrivenConstitutiveRelation}.  In that case:

\begin{itemize}
  \item the site stores the algorithmic internal variables;
  \item \texttt{compute\_response()} and \texttt{tangent()} query the relation
        with that explicit state;
  \item \texttt{update()} commits the site state through
        \texttt{relation.commit(alpha, eps)};
  \item \texttt{internal\_state()} exposes the site-level algorithmic memory,
        not the compatibility state embedded inside the relation.
\end{itemize}

This is the first place where the constitutive site becomes semantically more
important than the relation object.

\section{Interaction with \texttt{ConstitutiveIntegrator}}

The bridge integrator \texttt{EmbeddedRelationIntegrator} now detects the new
contract and prefers:

\begin{center}
\texttt{relation.commit(site.algorithmic\_state(), eps)}
\end{center}

over the legacy

\begin{center}
\texttt{relation.update(eps)}
\end{center}

when that explicit-state path exists.

Therefore the architectural hierarchy is now more coherent:

\begin{itemize}
  \item \texttt{PlasticityRelation} provides a local nonlinear map with an
        explicit algorithmic state argument;
  \item \texttt{MaterialInstance} owns the state of the constitutive site;
  \item \texttt{ConstitutiveIntegrator} drives the commit path.
\end{itemize}

\section{What the New Tests Prove}

The new regression tests verify three important facts:

\begin{enumerate}
  \item committing an external plasticity state evolves that external state;
  \item doing so does \emph{not} mutate the compatibility state owned by the
        naked relation;
  \item a \texttt{MaterialInstance} built from \texttt{PlasticityRelation} now
        stores the plastic internal variables at the site level, and subsequent
        stresses depend on that site state.
\end{enumerate}

That is the first concrete proof that behaviour is starting to move away from
law-owned memory and toward site-owned memory.

\section{Why the Cache Was Not Reused}

The embedded-state implementation of \texttt{PlasticityRelation} contained a
mutable cache keyed effectively by the strain.  That cache was valid only under
the implicit assumption that a given relation instance corresponded to a single
local state.

Once the same relation can be evaluated with arbitrary explicit states, that
assumption is no longer valid.  Reusing the old cache for the explicit-state
path would have been semantically unsafe.

The current explicit-state overloads therefore prioritise correctness over
micro-optimisation and compute directly from the supplied state.  If needed
later, the correct optimisation point is the constitutive site or the local
integrator, not a state-agnostic mutable cache inside the law.

\section{Current Boundary of the Refactor}

At the stage described in this chapter, only \texttt{PlasticityRelation} had
been migrated to the new explicit-state contract.  The following models still
kept their irreversible memory internally:

\begin{itemize}
  \item \texttt{MenegottoPintoSteel}
  \item \texttt{KentParkConcrete}
  \item \texttt{FiberSection}
\end{itemize}

This is intentional.  The architectural pattern is now proven on one important
continuum inelastic law before being propagated to uniaxial hysteretic and
section-reduction models.

That propagation is completed for the uniaxial steel and concrete laws in the
next chapter.

\section{Next Logical Step}

The next meaningful move is to propagate the same pattern to uniaxial
nonlinear laws, especially those used in reinforced-concrete fiber sections.
That will make the constitutive-site state explicit not only for continuum
plasticity, but also for steel hysteresis, concrete crushing/softening, and
eventually section-level nonlinear reductions.

Once that is in place, the path toward dynamic constitutive updates, finite
kinematics, nonlocal regularisation, and continuum-scale damage or fracture
becomes much clearer.
